\section{The non-vacuity threshold: protocol, tables and limits}
\label{sm:frontier}

Throughout this section we call $\mathrm{FR}^{\ast}$ the \emph{non-vacuity threshold of the certificate
of Theorem~\ref{thm:leakage}}: the sector fraction above which the leakage branch of the bound first
improves on the trivial bound $\lVert\phi\rVert^{2}(1+w_S)$. We do not call it a frontier of the
method, and it is not one. It is a property of the inequality chain that proves
Theorem~\ref{thm:leakage}, measured here for the first time, together with one structural fact about the determinant basis. Two measurements separate the two contributions. First, at
fixed system and fixed accuracy target the threshold sits far above the fraction the reconstruction
actually needs: at $L=8$, $\eta=0.18\,t$, a relative $L_1$ error of $0.05$---the acceptance criterion
used throughout this paper---is reached at $\mathrm{FR}\approx0.52$, while the certificate does not
reach $0.05$ until $\mathrm{FR}\approx0.98$, and the intervening shift of $\approx0.45$ in fraction is
the looseness of the proof. Second, at full captured weight $\Lambda_S=0$ requires $\operatorname{span}(S)\supseteq\mathcal K(H,\phi)$, whose coordinate support is every determinant the reflection about the seed site allows, at every size measured---a statement that the leakage is nonzero, not about its size, and not the source of the vacuity measured below, which is at $w_S<1$---and a Krylov-diagonal oracle ranking gives a \emph{worse} certificate than the Born ranking the device realizes at $L=6$, $\eta=0.18\,t$, at every fraction checked from $0.50$ on.
Section~\ref{sec:frontier:obstruction} measures both. What is a property of the method, and is new, is
the calibration of Sec.~\ref{sec:frontier:calib}: the crossing occurs at an essentially fixed level of
true error, so the crossing itself is worth computing even where the bound around it is not.

\paragraph{What is evaluated, in full.} We evaluate Theorem~\ref{thm:leakage} on the subspaces of Fig.~\ref{fig:scaling} at $L=6$, $8$, $10$
and $12$---sector dimensions $300$, $3920$, $52\,920$, $731\,808$---with the ranking protocol of
Sec.~\ref{sec:prices} ($K=18$, $\delta t=0.5$, $16$ substeps, $19$ snapshots), over fractions
$0.20$--$0.95$ and $\eta=0.10$--$0.25\,t$. Those subspaces are the $T\to\infty$ limit of that protocol
and not finite-shot draws, and every statement below inherits that limit: the marginal determinant at
the $L=12$ threshold carries Born probability $1$--$2\times10^{-9}$, so at the budgets of
Table~\ref{tab:shots} that subspace barely exists as a sampled object. Of the $204$ evaluations, $187$
are converged in Krylov depth, $6$ marginal and $11$ not converged, at a total cost of $2874$\,s with
a peak of $1.64$\,GiB. \emph{No row flagged not converged is quoted anywhere in this section}; the
tables mark the ones that are marginal.

To that grid we add twenty evaluations that are not on it. A grid point near a published fraction is
not the published subspace, and the difference is not cosmetic: this sweep reads $0.80$, $0.56$,
$0.40$ and $0.18$ against published fractions of $0.8200$, $0.5561$, $0.3625$ and $0.1700$. The rows
of Table~\ref{tab:published} are placed instead at the integer determinant count
$\lvert\mathcal S\rvert=\mathrm{round}(\mathrm{frac}\times\dim)$ that Fig.~\ref{fig:scaling}
reports---$246$, $2180$, $19\,183$, $124\,407$ and $824\,504$ determinants at $L=6$, $8$, $10$, $12$
and $14$---at the same four resolutions and under the same protocol, at depths $n_\ell=245$
(exhausting the subspace), $900$, $700$, $300$ and $250$, with orthogonality defect at most
$4.9\times10^{-15}$ and $\lvert w_{\mathcal K}-w_{\mathcal S}\rvert$ at most $6.4\times10^{-15}$. All
twenty pass the depth test, and all twenty clear the reference-drift rule of Sec.~\ref{app:hyp}, the
weakest---$L=14$ at $\eta=0.10\,t$, where the reference itself drifts by $1.9\times10^{-3}$ on
halving its depth---by a factor $51$. They cost $129$\,s of core time up to $L=12$ and $2238$\,s of
wall time at $L=14$, and they add the size the sweep never reached.

Three conventions are fixed before any number is read. Full reorthogonalization is used throughout,
without which the projector is not a projector (Sec.~\ref{app:hyp} gives the measured defect). What
is certified is the object the code forms---$\Lambda$ on
$\mathcal K=\mathrm{Krylov}(H_S,\phi_S,n_\ell)$ bounds the error
of $A_{\mathcal K}$ at that depth and of nothing else, the convenient inequality
$\Lambda_{\mathcal K}\ge\Lambda_S$ that would extend it to the dense $A_S$ being false ($1169$ of
$6000$ random triples violate it, minimum ratio $0.5258$), and $A_S$ is never formed. And depth is
certified per point by the internal indicator $\Lambda_{\rm in}/\Lambda_{\rm out}$ rather than
assumed~\cite{haydock1972,epperly2022}; Sec.~\ref{app:hyp} bounds what the depth cap does to
$\mathrm{FR}^{\ast}$.

Three hypotheses carry the bound, each a statement about the pipeline rather than about the theorem;
they are stated in full with the theorem in Sec.~\ref{sec:leakage}, and all three are checked against
this pipeline in Sec.~\ref{app:hyp}, along with the ranking, tie and support declarations of the cut.
It is (H3), the exactness of the reference, that enforces the rule that no measured error is quoted
within a factor $20$ of its own reference drift.
One asymmetry runs through every table: of the two terms of Eq.~\eqref{eq:twobudget} only the leakage
is a posteriori, while the weight term needs $\Psi_0$, and estimating $w_S$ against an
already-truncated $\Psi_0$ is biased optimistically (Sec.~\ref{sec:leakage}). What the tables report is
accordingly one object and not a family of them: a bound on
$\lVert A-A_{\mathcal K}\rVert_1/\lVert\phi\rVert^{2}$ for the reorthogonalized Krylov reconstruction
at a declared depth, against a reference whose drift is declared and enforced.

\paragraph{The two largest sizes.} The largest classically converged case carries the evidence. At $L=12$ and the published subspace
itself---$\lvert S\rvert=124\,407$ of $731\,808$ determinants---the captured weight is
$w_S=0.998622$, so Eq.~\eqref{eq:tau} asks for $\widehat\Lambda_S/\eta<\tau(w_S)=0.9625$ and the
measured value is $4.236$ at $\eta=0.18\,t$, putting the leakage branch at $8.544$ against a trivial
branch of $1.9986$. That is converged in depth: the depth indicator is
$\Lambda_{\rm in}/\Lambda_{\rm out}=2.8\times10^{-3}$ and between depths $250$ and $300$ $\Lambda_{\rm out}$ moves by $3.7\times10^{-7}$ relative and the bound by $1.4\times10^{-5}$, while at the neighbouring grid point $\mathrm{FR}=0.18$ the same
quantity was checked twice more---$\Lambda_{\rm out}$ changing by $7.5\times10^{-7}$ between depths
$250$ and $300$, and an independently written re-implementation moving it by $8.6\times10^{-5}$
between depths $100$ and $300$. In all sixteen deposited depth series $\Lambda$ at every depth from $100$ on lies at or above its deepest value---it is not monotone in depth in general, and rises by up to $2.2\%$ from one depth to the next below depth $100$ at $L=12$---so the bound at any depth from $100$ on, including the depth $150$ at which Fig.~\ref{fig:scaling} is drawn, is no smaller, so the vacuity is not an artifact of Lanczos truncation in either direction.
Nor is it peculiar to these subspaces: over the $606$
live subspaces of the pooled sweep of Sec.~\ref{app:counterex}, Fig.~\ref{fig:thm1iii-violation}, the
leakage branch is the smaller of the two on only $61$. \textbf{The reconstructions at the fractions this paper publishes---Fig.~\ref{fig:scaling} and the operating points of Secs.~\ref{sec:prices} and~\ref{sec:physics}---are therefore, at present, uncertified.} Fig.~\ref{fig:akwsampled}, drawn at $85\%$ of its sector, is the exception: there the certificate of the depth-$500$ Krylov reconstruction is non-vacuous in all sixteen channels, at a relative bound of $0.77$--$1.75$ against the trivial $2$, and loose, $350$--$900$ times the true error (\texttt{cert\_akw.json}).

One row is quoted with two caveats attached. $L=14$ is the size at which Fig.~\ref{fig:scaling}
reports its headline fraction of $0.08$, and the certificate is evaluated there at $\lvert S\rvert=824\,504$ of $10\,306\,296$ determinants---one determinant above the
$\mathrm{round}(\mathrm{frac}\times\dim)=824\,503$ of the deposited grid---with $w_S=0.998388$ and a
leakage branch of $17.21$, $11.18$, $9.21$ and $6.50$ against a trivial branch of $1.9984$. Ground
state, seed and order there are the repository's stored checkpoint---residual $2.7\times10^{-9}$, the
order verified to be a permutation of the sector---and the $K=18$ ranking protocol was \emph{not}
re-run at that size; it was re-run at $L=10$, where the two selections agree as sets to
$1.000000$. And the
depth indicator $\Lambda_{\rm in}/\Lambda_{\rm out}=0.0453$ at $\eta=0.10\,t$ is the loosest value
quoted anywhere in this section, comparable to the $4.7\times10^{-2}$ of the last row of
Table~\ref{tab:certified}. Neither caveat is repaired, and the row is reported because it makes the
verdict of this section worse rather than better.

\begin{table*}[tp]\centering\small
\caption{\textbf{Where the certificate stops being vacuous, and what the reconstruction actually
needs, measured on the same sweep.} Protocol of Sec.~\ref{sec:prices}, full reorthogonalization, depth
certified per point. Column~3 is the fraction Fig.~\ref{fig:scaling} reports; column~4 the control,
where this sweep---independent code, the normalization of Theorem~\ref{thm:leakage}, the whole line,
a deep reference---first reaches relative $L_1<0.05$, the parenthesized value being the log-linear
interpolation inside the bracket and not a measurement at three digits; column~5 is
$\mathrm{FR}^{\ast}$; columns~6--7 are where the certified value itself first falls below $0.50$ and
$0.20$. Every bracket is delimited by two converged evaluations unless marked: $^{\ast}$ means the
upper end is marginal in depth, \emph{n.c.}\ that the bracket rests only on evaluations failing the
depth test and is therefore not quoted. Columns~5--7 are statements about $A_{\mathcal K}$ at the
declared depth, not about the dense $A_S$. Columns~4 and~5 must be read together: $\mathrm{FR}^{\ast}$
alone falls with size and would suggest the certificate improves as the problem grows, whereas what
the reconstruction needs falls with it and the gap widens---a gap we do not quote as a ratio, both
columns being interpolations inside brackets $0.02$--$0.10$ wide.}\label{tab:thresholds}
\setlength{\tabcolsep}{4pt}
\begin{tabular}{lcccccc}
\toprule
$L$ & $\dim(N{+}1)$ & published & rel-$L_1<0.05$ & $\mathrm{FR}^{\ast}$ & cert $<0.50$ & cert $<0.20$\\
\midrule
\multicolumn{7}{l}{\emph{$\eta=0.18\,t$, the resolution of the momentum-resolved channel}}\\
$6$  & $300$     & $0.82$ & $[0.70,0.80]$ & $[0.90,0.92]$ & $>0.97$ & $>0.97$\\
$8$  & $3\,920$  & $0.56$ & $[0.50,0.56]$ & $[0.75,0.80]$ & $[0.90,0.92]$ & $>0.95$\\
$10$ & $52\,920$ & $0.36$ & $[0.30,0.40]$ & $[0.60,0.70]$ & $[0.80,0.85]$ & $[0.85,0.90]$\\
$12$ & $731\,808$& $0.18$ & $[0.14,0.18]$ & $[0.40,0.45]$ & $[0.60,0.70]^{\ast}$ & \emph{n.c.}\\
\midrule
\multicolumn{7}{l}{\emph{$\eta=0.15\,t$, the resolution at which Fig.~\ref{fig:scaling} is drawn}}\\
$6$  & $300$     & $0.82$ & $[0.80,0.85]$ $(0.816)$ & $[0.92,0.95]$ & $>0.97$ & $>0.97$\\
$8$  & $3\,920$  & $0.56$ & $[0.50,0.56]$ $(0.553)$ & $[0.80,0.82]$ & $[0.92,0.95]$ & $>0.95$\\
$10$ & $52\,920$ & $0.36$ & $[0.30,0.40]$ $(0.356)$ & $[0.60,0.70]$ & $[0.80,0.85]$ & $[0.85,0.90]$\\
$12$ & $731\,808$& $0.17$--$0.18$ & $[0.14,0.18]$ $(0.168)$ & $[0.45,0.50]$ & $[0.60,0.70]^{\ast}$ & \emph{n.c.}\\
\bottomrule
\end{tabular}
\end{table*}

One column of Table~\ref{tab:thresholds} is a positive control and has to be read before the rest.
Applying this paper's own acceptance criterion to the same sweep---independent code, the theorem's
normalization, the whole line rather than a window, a deep reference rather than one recomputed at the
subspaces' own depth---returns $0.816$, $0.553$, $0.356$ and $0.168$ at $\eta=0.15\,t$ for
$L=6$--$12$, against the published $0.82$, $0.56$, $0.36$ and $0.18$: agreement to $7\%$ at worst, and
to under $2\%$ at every size once each value is compared with the fraction the scan itself reports
rather than with the two-figure entry of the replication grid ($0.5561$, $0.3625$, $0.1700$ at
$L=8,10,12$). None of the three ingredients is shared. Figure~\ref{fig:scaling} is thereby reproduced
independently.

The last column of Table~\ref{tab:published} is that control taken one step further, and it is the
sharper of the two: instead of re-deriving the fraction it evaluates the relative $L_1$ error
\emph{at} the published subspace, at the resolution and against the threshold that define it
($\eta=0.15\,t$, rel-$L_1<0.05$). The five values---$0.0477$, $0.0478$, $0.0484$, $0.0485$ and
$0.0463$---all sit just below the threshold and are flat to within $5\%$ of one another across a
sector dimension running from $300$ to $10\,306\,296$, $L=14$ included. The criterion that
\emph{defines} the published fractions is therefore reproduced at the subspaces themselves, at every
size the figure plots above $L=4$, again with nothing shared. What fails in this section is not the
reconstruction and not that curve; it is the word
\emph{certificate} attached to the fractions at which the curve is read.

% [t] and not [tp], MEASURED 2026-09-19: with [tp] this table and
% tab:shotfrontier each qualified for a float-only page, and the queue spent a whole
% page on one small table with 60 %% of it blank.  Denying them 'p' costs nothing
% (no float is stuck, none is too large) and takes the manuscript 60 -> 59 pages.
\begin{table*}[t]\centering\small
\caption{\textbf{The threshold in the only currency a sampler has.} $T_{\rm marg}$ is the shots per snapshot at
which the marginal determinant of the cut appears once in expectation---the reciprocal of its Born
probability in the accumulated distribution of the $19$ snapshots, from the protocol's own cut
weights; covering all of $S$ costs a further factor $\log\lvert S\rvert$. \emph{It is a different
quantity from the measured budgets $T$ of Table~\ref{tab:shots}}, which are totals over the $19$
snapshots at which the finite-shot reconstruction reaches relative $L_1<0.05$; the two carry different
names for that reason, they are not interchangeable, and only ratios within this table are used.
Column~2 uses the grid point nearest the published fraction: $0.80$ against a published $0.82$ at
$L=6$, and $0.40$ against $0.36$ at $L=10$, where the published fraction is the worse of the two.
Factors are from unrounded budgets and need not reproduce from the two-figure entries. The marginal Born probability at
the $L=12$ threshold is $2.2\times10^{-9}$ at $\mathrm{FR}=0.40$, $1.2\times10^{-9}$ at
$0.45$.}\label{tab:shotfrontier}
\setlength{\tabcolsep}{6pt}
\begin{tabular}{lccc}
\toprule
$L$ & published fraction $\to T_{\rm marg}$ & $\mathrm{FR}^{\ast}\to T_{\rm marg}$ & factor\\
\midrule
$6$  & $0.80\to2.2\times10^{4}$ & $[0.90,0.92]\to3.1$--$3.6\times10^{4}$ & $1.4$--$1.6$\\
$8$  & $0.56\to2.1\times10^{5}$ & $[0.75,0.80]\to1.6$--$2.1\times10^{6}$ & $7.5$--$10$\\
$10$ & $0.40\to5.2\times10^{6}$ & $[0.60,0.70]\to3.4\times10^{7}$--$1.1\times10^{8}$ & $6.7$--$21$\\
$12$ & $0.18\to2.3\times10^{7}$ & $[0.40,0.45]\to4.7$--$8.1\times10^{8}$ & $20$--$35$\\
\bottomrule
\end{tabular}
\end{table*}

In shots the distance between those two columns is a budget and not a rhetorical gap
(Table~\ref{tab:shotfrontier}): reaching the non-vacuous band multiplies the shots per snapshot needed
for the marginal determinant by $7.5$--$10$ at $L=8$ and by $20$--$35$ at $L=12$, and only by
$1.4$--$1.6$ at $L=6$, where the sector is small. Driving the certified value down to $0.10$ costs more
again---a further factor $241$ at $L=8$ (fraction $0.95$) and $634$ at $L=10$ (fraction $0.90$), the
corresponding $L=12$ fraction resting only on evaluations that failed the depth test, so that no budget
is quoted for it. This section reports the measured regime of validity: four sizes, four resolutions, $204$ evaluations, every
bracket delimited by converged rows, and a shot budget attached to each.

\paragraph{The strongest certified statements.} Table~\ref{tab:certified} collects what the sweep does certify, taking at each size the strongest
statement whose bracket is converged---at $L=12$, that $60\%$ of a $731\,808$-dimensional sector gives
relative $L_1\le0.717$, tightening to $\le0.218$ at $70\%$ and $\eta=0.25\,t$. Each row is a statement
about $A_{\mathcal K}$ at the declared depth, and none of them uses the exact spectrum.

% [t] and not [tp], MEASURED 2026-09-19: with [tp] this table and
% tab:shotfrontier each qualified for a float-only page, and the queue spent a whole
% page on one small table with 60 %% of it blank.  Denying them 'p' costs nothing
% (no float is stuck, none is too large) and takes the manuscript 60 -> 59 pages.
\begin{table*}[t]\centering\small
\caption{\textbf{The strongest certified statements available from converged data.} Each row is the
tightest certified value at that size whose bracket is delimited by evaluations passing the depth
test; rows that would be stronger but rest on evaluations failing it are not listed. ``True error'' is
the measured relative $L_1$ at the same point, in the normalization of Theorem~\ref{thm:leakage},
given only where it exceeds its own reference drift by more than a factor $20$ (Sec.~\ref{app:hyp},
H3); at $L=10$ it does not, and is withheld rather than rounded. At $L=6$ the recursion runs to
$n_\ell=\lvert S\rvert$, so the evaluation is exact and the depth indicator is not defined. The gap
between the certified value and the true error is three to four decades: the object is a certificate
and not an error estimate, and Sec.~\ref{sec:frontier:calib} states what it is good
for.}\label{tab:certified}
\setlength{\tabcolsep}{5pt}
\begin{tabular}{lccccc}
\toprule
$L$ & $\eta/t$ & fraction & certified rel-$L_1\le$ & true rel-$L_1$ & $\Lambda_{\rm in}/\Lambda_{\rm out}$\\
\midrule
$6$  & $0.18$ & $291/300=0.97$            & $0.502$ & $3.94\times10^{-4}$ & exact ($n_\ell=\lvert S\rvert$)\\
$8$  & $0.18$ & $3\,724/3\,920=0.95$      & $0.219$ & $4.33\times10^{-5}$ & $8\times10^{-4}$\\
$10$ & $0.18$ & $47\,628/52\,920=0.90$    & $0.136$ & withheld (H3)       & $1.6\times10^{-3}$\\
$12$ & $0.18$ & $439\,085/731\,808=0.60$  & $0.717$ & $1.89\times10^{-4}$ & $7.0\times10^{-3}$\\
$12$ & $0.15$ & $0.60$                    & $0.878$ & $2.66\times10^{-4}$ & $1.4\times10^{-2}$\\
$12$ & $0.25$ & $512\,266/731\,808=0.70$  & $0.218$ & $1.99\times10^{-5}$ & $4.7\times10^{-2}$\\
\bottomrule
\end{tabular}
\end{table*}

The concession that belongs with that table is one a referee would make, so we make it first and in
full. At the acceptance criterion this paper uses everywhere else---relative $L_1<0.05$---the
certificate requires at least $0.98$ of the sector at $L=8$, measured at integer subspace size rather
than interpolated: the certified value is $0.084$ at $3802$ of $3920$ determinants, $0.060$ at $3830$
and $0.039$ at $3855$, so the crossing lies in $[0.977,0.983]$; on our grid it is not reached at all at
$L=6$, $L=10$ or $L=12$. At such fractions diagonalizing the whole sector is already what one would do
instead, and a certificate informative only there does not certify the regime in which sampling is of
interest. We state it without qualification: at the accuracy this paper accepts elsewhere,
Theorem~\ref{thm:leakage} does not certify the operating points of Fig.~\ref{fig:scaling}, and the shot
cost of the regime where it would exceeds the saving that figure reports.

What the table does certify at a fraction a sampler can plausibly reach is weaker, and it is not
empty. The trivial bound caps the relative error at $1+w_S\le2$ for \emph{any} subspace, a
reconstruction of pure noise included; at $L=12$ and $60\%$ of a $731\,808$-dimensional sector the
certified $0.717$ sits a factor $2.8$ below that cap, and at $70\%$ and $\eta=0.25\,t$ a factor $9.2$
below it. Those are exclusions of gross failure made from the subspace alone, at a dimension where the
exact spectrum is expensive and, at the sizes this method is aimed at, unavailable.

\paragraph{The calibration and its limits, in full.} One positive result survives the sweep, and it is the reason we report the sweep. Across all sixteen
$(L,\eta)$ cells the crossing of the trivial bound occurs at a true relative $L_1$ error of
$1.2\times10^{-3}$ to $7.3\times10^{-3}$, median $2.8\times10^{-3}$: a spread of $5.9$ over a table
whose true errors span five decades, and in every cell at least $47$ times the drift of our own
reference. $\widehat\Lambda_S(\eta)/\eta$ is therefore a calibrated a posteriori risk tracker---its
crossing of the trivial bound locates a true error of ${\approx}3\times10^{-3}$ to within a factor of
three over $L=6$--$12$ and $\eta=0.10$--$0.25\,t$, and it is computed from $H$, $S$, $\phi$
and $\eta$ alone, without the answer. That factor of three is a drift and not scatter, and the
deposited grid settles it in a few lines: the sixteen crossing errors are ordered without a single
exception, falling with $L$ at every resolution---$5.65$, $3.50$, $2.52$ and $1.60\times10^{-3}$ at
$\eta=0.18\,t$ for $L=6$--$12$, a factor $2.8$--$3.7$---and rising with $\eta$ at every size by
$1.6$--$2.1$, the product of the two being the observed $5.9$. The drift runs the conservative way in
the variable one would extrapolate: the larger the system, the smaller the true error at which the
certificate first says anything, so reading the crossing as ${\approx}3\times10^{-3}$ beyond the sizes
measured overstates the error rather than understating it. In $\eta$, which the user chooses and which
lies inside the measured range, it does not. The objection to anticipate is that this is
Fig.~\ref{fig:scaling} relabelled. It is not, and the difference is the content: that figure is built
knowing the exact spectrum and this quantity is computed without it. That the two track each other to
within a factor $5.9$ across five decades of error is not a redundancy---it is the calibration, and it
is what makes the quantity usable where the second curve cannot be drawn.

Four limits come with it. The crossing is one point per cell, so the tracker locates one error
level and does not estimate the error away from it: at a \emph{fixed} subspace the looseness of the
bound against the measured error is constant to within $\pm17\%$ ($L=6$) and $\pm25\%$ ($L=8$) over a
decade in $\eta$, but across subspaces it varies by a factor $51$, and at the threshold itself it is
$3\times10^{2}$ to $2.4\times10^{3}$---the object is a certificate, not an error estimate, and calling
it the latter would be refutable in one line. Only one of its two terms is a posteriori: the comparison is
$\widehat\Lambda_S/\eta$ against $\tau(w_S)$, and $\tau$ needs $w_S$, which a user does not have.
That dependence can be discharged conservatively, and doing so is what makes the decision itself a
posteriori. Equation~\eqref{eq:tau} is decreasing in $1-w_S$ for $w_S\ge0.066$, with $\tau(1)=1$, and positive for $w_S>0.296$, so any assumed floor on the captured weight above $0.296$ yields a usable threshold: at the $L=12$ published operating point
$w_S=0.998622$ and $\tau=0.9625$, and every subspace with $w_S\ge0.995$ has $\tau\ge0.9280$. A user
willing to assume only that the sampler captured $99.5\%$ of the probe's Born weight can therefore
run the comparison as $\widehat\Lambda_S(\eta)/\eta<0.928$ and use no other ground-state
information, at the cost of at most $7\%$ in the threshold. Where not even that floor can be
assumed, or where $w_S$ is estimated against an already-truncated $\Psi_0$---an estimate biased
optimistically---the test is not conservative and the bound is not a posteriori.
And sixteen cells over four sizes and four resolutions is the whole of the evidence: no exponent is
fitted to it and no extrapolation to $L=14$ is made, the local log--log slopes of
$\mathrm{FR}^{\ast}(L)$ already varying by a factor $4.1$ across the four sizes measured.

The fourth limit is the one a reader holding the deposited grid will test first, and it is a real
restriction on what the crossing is worth. A sampler already holds the captured Born weight, so the
question is whether a stopping rule indexed on $1-w_S$ alone locates the same error level, in which
case the crossing would be that weight relabelled. It is not: over the sixteen cells, $1-w_S$ at the
crossing ranges over a factor $45$ against the true error's $5.9$, and still over $9.4$ against $3.8$
once the extreme cell at each end is dropped. A rule with constants fitted to the grid, however, comes
close enough. Take $1-w_S<c\,\eta^{p}$, fit $c$ and $p$ on some of the four sizes
with its operating point matched to the certificate's, evaluate it on the sizes held out, and repeat
over all fourteen train/test splits: out of sample it localizes the crossing error about as tightly as
the parameter-free certificate does, and more tightly on eleven of the fourteen (\texttt{null\_ws7.py}).
We therefore claim no superiority for the certificate as a tracker of the error against the captured
weight; out of sample it has none. The difference lies elsewhere. The fitted rule fires at $1-w_S\sim10^{-5}$--$10^{-4}$, one to two orders of magnitude finer than the floor
$1-w_S\le5\times10^{-3}$ at which the threshold above still works, and it asks that precision of the one
quantity this section has already said a user does not have and, estimated against an already-truncated
$\Psi_0$, has optimistically; strip its second parameter and an $\eta$-blind $1-w_S<c$ is the looser of the two on thirteen of the fourteen splits. Theorem~\ref{thm:leakage}, for its part, is an inequality with a proof:
it holds at a size no grid has seen, and its decision survives replacing $w_S$ by a floor. The fitted
rule is a two-constant regularity with no guarantee outside its fitted range---the only place either
would be used.

Within those limits the procedure is operational and identical at every size we measured: assemble
$\widehat\Lambda_S(\eta)/\eta$ from the Ritz vectors already in hand and compare it with $\tau(w_S)$.
Above $\tau$ the certificate says nothing, which is the regime Fig.~\ref{fig:scaling} operates in; at
$\tau$ the true relative error is ${\approx}3\times10^{-3}$ within a factor of three; below $\tau$ it
is smaller than that, by an amount the bound states loosely and the calibration does not sharpen. That
is the whole content of the certificate, and it is a number a user can compute before knowing whether
the reconstruction is right.

\paragraph{Where the threshold comes from, in full.} Where does the ${\approx}0.45$ of sector fraction between what the reconstruction needs and what the
certificate certifies go? We measured the displacement rung by rung along the chain of inequalities
that proves Theorem~\ref{thm:leakage}, at $L=8$, and the answer removes the one repair that looked
available: the pointwise step from a scalar distance to a distance in the Hilbert space accounts for
$44\%$ of it, the operator-norm step $\lVert R(z)\rVert\le1/\eta$ applied to a source that depends on
$\omega$ for $38\%$, the Feshbach/Minkowski split for $11\%$, and the Cauchy--Schwarz step in $\omega$
for $4.6\%$ (medians over $15$ of $18$ cells; Sec.~\ref{app:hyp}). The working hypothesis that
Cauchy--Schwarz is what places the threshold is therefore refuted by its own measurement: its total
margin for tightening is a factor $1.05$ typically and $3.29$ at the worst point measured, which moves
the threshold by a fraction of one grid interval.

The second measurement concerns the subspace itself, and it is the one that turns part of the result into a statement about the problem. $\Lambda_S=0$ holds if and only if $\operatorname{span}(S)$ contains $\mathcal K(H,\phi_S)$, the Krylov space of the retained part of the probe (Sec.~\ref{sec:leakage}); for a subspace that holds the whole probe, $w_S=1$, that space is $\mathcal K(H,\phi)$ itself, so the coordinate support of $\mathcal K(H,\phi)$ is a floor under every coordinate subspace that holds the whole probe without leaking. Measured, it is $296$ of $300$ determinants at $L=6$, where $\dim\mathcal K=65$, and $3902$ of $3920$ at $L=8$, where $\dim\mathcal K=596$--$663$. (It is the \emph{coordinate support} of $\mathcal K$ that is measured, not $\mathcal K$ itself, whose dimension is three orders of magnitude smaller; containment, not spanning, is what $\Lambda_S=0$ requires.) The determinants missing are not an accident of threshold. The reflection $j\to-j$ about the seed site commutes with $H$, and $\phi$ is an eigenvector of it, so every vector of $\mathcal K$ vanishes exactly on the reflection-fixed determinants of the opposite parity: $4$, $18$, $48$, $200$ and $600$ at $L=6$, $8$, $10$, $12$ and $14$. Every other determinant is in the support. At $L=10$ the same
construction---a Lanczos recursion on the whole sector to depth $700$ with full reorthogonalization,
orthogonality defect $1.6\times10^{-15}$---leaves no row exactly zero in floating point, although $48$ of them vanish in exact arithmetic: round-off is not a support.

Beyond $L=10$ that basis cannot be stored, and we use a lower bound instead: every $H^{k}\phi$ lies in
$\mathcal K$, so the union over $k$ of the coordinate support of $H^{k}\phi/\lVert H^{k}\phi\rVert$ is
contained in the support of $\mathcal K$. A count of this kind is meaningless without its threshold,
so both are given. Projected onto the reflection sector of $\phi$ at every power (\texttt{z\_suppK\_sym.py}), so that round-off cannot populate the forbidden sector, the union over $k\le36$ is every symmetry-allowed determinant at every size: $296$, $3902$, $52\,872$, $731\,608$ and $10\,305\,696$ at $L=6$ to $14$, at an amplitude threshold of $10^{-12}$, and the same at $10^{-10}$ except at $L=14$, where $232$ genuine amplitudes lie between the two thresholds. Unprojected, the power iteration amplifies round-off in the forbidden sector until it crosses either threshold and reports the entire sector from $L=8$ on (\texttt{z\_suppK\_power.py}); that count is not a support.

At every size measured, therefore, a coordinate subspace that holds the whole probe leaks unless it contains every symmetry-allowed determinant---the sector minus its forbidden determinants contains $\mathcal K$ and does not leak---so at full captured weight no choice of determinants short of all of them makes $\Lambda_S$ vanish, and that is a property of the basis in which the device measures, not of the bound and not of the ranking. It is a statement that $\Lambda_S\neq0$, not about its size, and the vacuity at the published fractions is not deduced from it: that is the measurement of Table~\ref{tab:published}. Below full captured weight the space that decides is $\mathcal K(H,\phi_S)$, which changes with $S$, and nothing structural is claimed there. Nor is the ranking the weak link: replacing the Born ordering by a Krylov-diagonal oracle---determinants ordered by the diagonal of the projector onto $\mathcal K(H,\phi)$, an ordering no sampler could produce---gives a \emph{worse} certificate at $L=6$, $\eta=0.18\,t$, at every fraction checked from $0.50$ on, $12.97$ against $4.27$ at $0.85$ and $5.82$ against $0.911$ at $0.95$; only at $0.30$, where both are far above the trivial value, is the oracle marginally better, $21.76$ against $24.52$.

At fixed size the displacement is the looseness of the proof, and part of the looseness has a structural reason: the leakage branch is of first order in the boundary coupling, the error it bounds of second. For a subspace that holds the whole probe, eliminating $Q$ exactly replaces $H_P$ by $H_P+PHQ(z-QHQ)^{-1}QHP$, and the two sum rules then give $\lVert A-A_P\rVert_1\le\lVert\phi\rVert^{2}\lVert QHP\rVert^{2}/\eta^{2}$. Sharper, with $\phi_Q=0$ Eq.~\eqref{eq:split} and Cauchy--Schwarz in $\omega$ give $\lVert A-A_P\rVert_1\le\Lambda_P(\eta)\tilde\Lambda_P(\eta)/\eta^{2}$, where $\tilde\Lambda_P$ is Eq.~\eqref{eq:lambda} with the energy difference in its kernel reversed in sign, equal to $\Lambda_P$ for a real Hamiltonian and probe. The ratio of the leakage branch, $2\lVert\phi\rVert\Lambda_P/\eta$, to the error is therefore at least $2\eta\lVert\phi\rVert/\tilde\Lambda_P$, and grows without limit as the subspace approaches invariance; at full captured weight the product $\Lambda_P\tilde\Lambda_P/\eta^{2}$ is itself an a posteriori certificate of second order. Every published fraction lies below full captured weight, and whether a certificate of second order exists there is open; we have not constructed one. Section~\ref{sec:nogo} closes the picture
from the other side: nothing cheaper than the crossing of Sec.~\ref{sec:frontier:calib}---no invariant
of the state, computable before the subspace exists---can be put in its place.

One route out of the determinant basis has been measured, and it closes this section as an outlook and
not as a result, because the same measurement shows where it stops. The qualifier \emph{in the
determinant basis} above is load-bearing, and it is measurable. An orbital rotation is a Givens network
and the circuit of Sec.~\ref{sec:method} already contains them, so the measurement basis is a choice;
for the momentum-resolved probe of Fig.~\ref{fig:akwsampled} it is a forced one, since $H$ commutes with
lattice translation and $\phi_k$ is an eigenvector of it, so that $\mathcal K(H,\phi_k)$ lies inside a
single momentum block---rigorously, at every $L$---and that block is exactly $\dim(N{+}1)/L$
determinants. At $L=6$ and $L=8$ the Krylov support fills nearly all of the block, and the certificate becomes non-vacuous only close to the full block: at $48$--$50$ of its $50$ determinants at $L=6$ and at $455$--$483$ of its $490$ at $L=8$, over the momenta of the addition branch (\texttt{z\_kblock.py}), a factor $5$--$7$ fewer than the determinant basis requires; on the block itself $\Lambda_S$ vanishes identically and the bound is zero rather than vacuous. The gain, though, is $1/L$ against an exponential: as a fraction of
the fraction Fig.~\ref{fig:scaling} publishes the block is $0.20$ at $L=6$ and then $0.22$, $0.28$,
$0.49$ and $0.89$ at $L=14$, the two crossing just past $L=14$. It is also a laptop Lanczos---$60\,984$
determinants at $L=12$, $736\,164$ at $L=14$---so the regime it certifies is one that momentum-resolved
exact diagonalization has enumerated since the 1970s. Three things are unmeasured here and would have to
be measured before more is said: whether the support fills the block at $L\ge10$, where the threshold
then sits, and what the ${\approx}L(L-1)$ extra two-qubit gates do to the sampled distribution.
